% Compile with XeLaTeX:  xelatex resume_bullets.tex
% Required fonts: any system CJK font works (macOS: PingFang SC / STSong)
%
\documentclass[11pt, a4paper]{article}
\usepackage[UTF8]{ctex}
\usepackage[T1]{fontenc}
\usepackage{geometry}
\usepackage{enumitem}
\usepackage{tabularx}
\usepackage{booktabs}
\usepackage{xcolor}
\usepackage{hyperref}
\usepackage{parskip}

\geometry{top=2cm, bottom=2cm, left=2.2cm, right=2.2cm}

% ── colours ──────────────────────────────────────────────────────────────────
\definecolor{headgray}{RGB}{60,60,60}
\definecolor{accentblue}{RGB}{30,100,180}
\definecolor{lightbg}{RGB}{245,247,250}

% ── helpers ──────────────────────────────────────────────────────────────────
\newcommand{\projhead}[2]{%
  \noindent{\large\bfseries\color{headgray} #1}%
  \hfill{\small\color{gray} #2}\par\smallskip
}
\newcommand{\versionlabel}[1]{%
  \medskip\noindent{\bfseries\color{accentblue} #1}\par\smallskip
}
\newcommand{\langlabel}[1]{%
  \noindent{\small\bfseries\color{gray} #1}\par
}

\setlist[itemize]{leftmargin=1.4em, itemsep=3pt, topsep=2pt}

% ─────────────────────────────────────────────────────────────────────────────
\begin{document}
\pagestyle{empty}

% ════════════════════════════════════════════════════════════════════════════
\projhead{Multimodal Fashion Image Retrieval \quad 多模态时尚图像检索}{%
  Python · PyTorch · CLIP · HuggingFace \quad CS5330, Northeastern University}
\vspace{2pt}
\hrule
\vspace{8pt}

% ════════════════════════════════════════════════════════════════════════════
\versionlabel{版本一：纯个人贡献（推荐）\quad Version 1 — Own Work Only (Recommended)}

\bigskip

% ── side-by-side ─────────────────────────────────────────────────────────────
\begin{tabularx}{\textwidth}{X X}
\toprule
\textbf{中文} & \textbf{English} \\
\midrule

%  bullet 1
将 FashionCLIP（时尚领域微调 ViT）集成至基于 CLIP 的组合式图像检索流水线，以领域感知
嵌入替代零样本文本-图像匹配
&
Integrated FashionCLIP (fashion-domain fine-tuned ViT) into a CLIP-based composed image
retrieval pipeline, replacing zero-shot text-image matching with domain-aware embeddings
\\[6pt]

%  bullet 2
发现并修复关键 gallery 构建 bug（18,136 $\to$ $\approx$5,000 验证集图像），该 bug
导致 Recall@10 被压制约 3$\times$；通过正确应用 CLIP 的 \texttt{visual\_projection}
与 \texttt{text\_projection} 层修复嵌入空间错配问题
&
Diagnosed and fixed a critical gallery construction bug (18,136 $\to$ \textasciitilde5,000
val-split images) suppressing Recall@10 by $\sim$3$\times$; corrected embedding space
mismatch by applying CLIP's \texttt{visual\_projection} and \texttt{text\_projection} layers
\\[6pt]

%  bullet 3
在三个 Fashion-IQ 类别（dress / shirt / toptee）上系统评估纯文本、纯图像及加权融合
（$\alpha$=0.7）策略；FashionCLIP 融合在 toptee 类别达到 R@10 = 33.0\%，对比基线
CLIP 的 21.1\%（相对提升 +57\%）
&
Evaluated text-only, image-only, and weighted fusion ($\alpha$=0.7) across three Fashion-IQ
categories; FashionCLIP fusion achieved R@10 = 33.0\% on toptee vs.\ 21.1\% for baseline
CLIP (+57\% relative improvement)
\\[6pt]

%  bullet 4  ← update this number after 25-epoch run finishes
基于 16,800 个训练三元组，以 InfoNCE 批对比损失训练残差 MLP Combiner（1.6M 参数，
FashionCLIP 主干冻结）；三类别平均 R@10 从线性融合的 29.3\% 提升至 \textbf{33.65\%}
（相对 +14.8\%），dress 类别相对提升最显著（+25.5\%，26.97\% $\to$ 33.85\%）
&
Trained a residual MLP combiner (1.6M params, frozen FashionCLIP backbone) with InfoNCE
batch contrastive loss on 16,800 triplets; avg R@10 improved from 29.3\% (linear fusion)
to \textbf{33.65\%} (+14.8\% relative), with +25.5\% relative gain on the dress category
(26.97\% $\to$ 33.85\%)
\\
\bottomrule
\end{tabularx}

\vspace{14pt}

% ════════════════════════════════════════════════════════════════════════════
\versionlabel{版本二：含队友部分（项目整体范围）\quad Version 2 — Full Team Scope}

\medskip
\noindent\textcolor{gray}{\small
  说明：第 1--4 条与版本一相同（个人贡献）；第 5--6 条描述队友负责的模块，
  因尚无量化结果建议措辞为"设计"或"探索"，不写具体数字。}

\bigskip

\begin{tabularx}{\textwidth}{X X}
\toprule
\textbf{中文} & \textbf{English} \\
\midrule

%  bullets 1-4  (same as above, condensed)
\multicolumn{2}{l}{\textit{\color{gray}（第 1--4 条与版本一相同，此处省略）}} \\[6pt]

%  bullet 5  — TIRG teammate
（团队）探索 TIRG（文本-图像残差门控）作为 Combiner 的对照方案，通过可学习门控
机制将修改文本融入参考图像嵌入，实现更精细的查询表示
&
(Team) Explored TIRG (Text-Image Residual Gating) as a contrast to the MLP combiner,
incorporating modification text into the reference image embedding via a learned gating
mechanism for finer-grained query representation
\\[6pt]

%  bullet 6  — Re-ranking teammate
（团队）设计两阶段重排序 pipeline：以 FashionCLIP 粗检索 top-50 候选，再用轻量
交叉注意力模型对候选集重新打分排序，在不增加检索规模的前提下提升精排精度
&
(Team) Designed a two-stage re-ranking pipeline: coarse retrieval of top-50 candidates
via FashionCLIP, followed by a lightweight cross-attention re-ranker to refine ordering
without expanding the retrieval pool
\\
\bottomrule
\end{tabularx}

\vspace{10pt}
\noindent\textcolor{accentblue}{\small\bfseries 建议：}
\textcolor{gray}{\small
版本二第 5--6 条无量化数字，在简历中会比前四条弱。
若队友后续产出了结果，补上 R@10 数字即可；否则建议使用版本一，更干净有力。}

\end{document}
